\chapter*{Preface}
\addcontentsline{toc}{chapter}{Preface}
\markboth{Preface}{Preface}

Every computer science degree teaches the same shape of thing badly, or not at all: the
layer that sits between the code we write and the machine that runs it. We learn
algorithms, data structures, operating system theory, sometimes even how a compiler works
internally, and then we sit down at a real machine, in a real terminal, and none of it
tells us why \texttt{python} with no arguments starts an interactive prompt, what actually
happens when we type \texttt{git push}, or why our carefully installed packages vanish the
moment we open a different Python environment.

This book exists because that gap is real, it is common, and it is entirely fixable. It
just needs someone to actually walk through it with us instead of assuming we will pick it
up by osmosis from Stack Overflow and trial and error.

We wrote it a chapter at a time, the same way it should be read: never faster than
understanding allows. Each chapter grew out of an actual session: a real command, typed
into a real terminal, on a real machine that already had years of accumulated mess sitting
on it. We did not invent clean examples for this book. We used what was actually there:
two Python installations nobody remembered installing, three hundred globally installed
packages, a \texttt{git} workflow that was three memorized commands deep and zero commits
deep in understanding. If the machine in front of you looks like that too, good: that is
exactly who this book is for.

A few things about how we have chosen to write this.

We use ``we'' throughout, not ``you,'' because we do not think of this as a lecture. It is
closer to two people sitting at the same terminal, one of whom has already been down this
particular road and knows where the potholes are. When we say the prompt changed, or the
command failed in a particular way, we mean it: we are looking at the same output
together.

We do not skip the history. Every tool in this book, the shell, virtual environments,
Docker, Git itself, exists because something else broke first, and understanding what
broke is most of what makes the current design stop looking arbitrary. A rule with no
history behind it is just something to memorize and forget. A rule with its history
attached is something we can reconstruct ourselves, years later, even after the exact
syntax has slipped away.

We assume the degree, not the practice. This book is for computer science enthusiasts who
have done the courses: we know what a process is, how virtual memory works, what a compiler
does, what TCP promises, what ACID means. We will not re-teach any of that. When a chapter
leans on a course concept, it names it in a sentence and moves on to what the courses
skipped: what that concept looks like on a real machine, which tool shows it to us, and how
it breaks. What we do not assume is practical fluency, because that is the gap this book
fills, and it is usually invisible to the people who have it. It hides behind a Run button,
a pre-built cloud notebook, or a friend who set things up once, years ago, and it hides so
well that an entire degree can be finished without ever noticing it is there. So the first
time any command, file, or tool appears, we introduce it properly.

No two chapters look quite alike, and that is deliberate rather than sloppy. A chapter
that is mostly a mental model reads as prose, with a margin note where a definition or a
piece of history earns one. A chapter that lives in the terminal shows the real command
and the real output. A chapter that only makes sense side by side with three other
ecosystems becomes a table instead of a wall of text pretending to be one. The shape
follows what is actually being taught. What stays constant, chapter to chapter, is the
questions at the end: real ones, heard in real conversation, not invented to fill a
section.

\theading{How to use this book}{questions, not commands}

There is nothing here to memorize. The commands will change; the questions will not. For
any tool, in any chapter, we want to be able to answer: what is this tool actually doing,
and at which layer? What does it talk to? Where does its configuration come from? What
happens when it fails, and how would we know? How can we look inside it while it runs? Each
part ends with a trace: one real operation followed through every layer we have covered so
far, so the parts add up to one machine rather than a list of tools. Each chapter also
carries a tag in the margin, core, practical, deep dive, or specialized, so that we can
choose our own path through it. And the companion repository is the laboratory: the book
explains, the exercises make us do it, break it, and repair it ourselves.

If this book leaves us with working commands but no mental model of why they work, it has
failed at the one thing it set out to do. Everything else, the terminal fluency, the
\texttt{git} confidence, the Docker comfort, is downstream of that.

Let's begin.
